\subsection{Sprint 3}
Após a reunião com o stakeholder, o time foi liberado para começar os preparativos para a próxima sprint. Além de abordar os débitos técnicos que surgiram após a segunda sprint, nos comprometemos a desenvolver User Stories relacionadas a consultas, validações e emissões de certificados. Todas essas tarefas estavam mais relacionadas ao backend/frontend do que à própria blockchain. Vale ressaltar que nossa API no Ethereum ainda não estava completamente funcional, representando um débito técnico a ser resolvido.

Essa seria uma sprint muito focada no backend, então eu teria bastante trabalho pela frente. Para me preparar, dediquei a primeira semana quase exclusivamente a me aprofundar mais sobre aplicações SpringBoot. Fiz isso porque na sprint 1 senti que me faltava conhecimento sobre o framework, então tirei um tempo para estudá-lo. Foi uma semana muito produtiva, pois consegui aprender e reaprender de forma mais eficaz conceitos de backend que antes não estavam claros para mim. Essa dedicação me preparou bastante para as tarefas que estavam por vir e certamente me ajudou a realizá-las com mais facilidade e confiança.

A primeira tarefa que peguei para esta sprint foi listar todos os certificados de uma empresa. Essa tarefa, em específico, foi algo novo para mim, pois exigia trabalhar com aspectos como paginação, headers e autenticação em uma única task. Para conseguir realizá-la, pedi ajuda aos colegas de AGES IV, que foram muito compreensivos com minhas habilidades e me orientaram através de uma aula. Depois dessa orientação, senti que seria capaz de completar a tarefa, e foi exatamente o que aconteceu. Durante a execução, pude aplicar o que aprendi sobre SpringBoot e backend. Comecei configurando os endpoints e as requisições de paginação. Em seguida, trabalhei na inclusão de headers necessários para autenticação e gerenciamento de sessões. A integração com o banco de dados também foi um desafio interessante, mas consegui mapear corretamente os dados e garantir que a consulta fosse eficiente e após algumas horas de trabalho e testes, finalizei a tarefa com sucesso.

\sloppy

Para a segunda tarefa, resolvi trabalhar em uma task relacionada a "deletar" funcionários de uma empresa. Basicamente, a ideia por trás dessa operação era a deleção lógica de um funcionário. Em vez de realmente deletar um funcionário do banco de dados, implementaríamos um atributo, como "conta ativa", que indicaria se a conta do funcionário estava ativa ou desativada. Optamos por essa abordagem porque os certificados, na modelagem do banco de dados, dependiam do funcionário. Isso significava que, se realmente deletássemos o usuário, todos os certificados relacionados a ele também seriam deletados. Para contornar esse problema, utilizamos a ideia de deleção lógica. A task em si foi bem tranquila de realizar, uma vez que eu havia aprendido muito com meus colegas de AGES IV e com a semana de estudos que fiz no início da sprint. A tarefa em si envolveu a criação do endpoint para a operação, a configuração da lógica de autenticação, a implementação da lógica da tarefa no Service e, por fim, a conexão ao banco de dados para alterar o estado do funcionário.

Apesar da Sprint estar sendo muito produtiva, acabei me deparando com um grande problema. Meus Merge Requests (MRs) das tarefas que desenvolvi durante a sprint apresentaram um sério problema. Minhas branches, por algum motivo, se juntaram com a develop, causando muitos conflitos e impedindo a aceitação dos MRs. Isso me desanimou bastante, pois o trabalho para tentar resolver esses conflitos foi imenso e ainda assim não foi o suficiente. Depositei minhas esperanças nos colegas de AGES superiore para me auxiliar nessa situação, e, no final, conseguimos resolver o problema. Mesmo assim, essa experiência foi bastante desanimadora para mim.

A Sprint teve seu término com nossa penúltima reunião com o stakeholder. Desta vez, as aulas presenciais já haviam retornado, mas Edimar não conseguiu comparecer presencialmente à reunião. Em geral, a sprint foi muito produtiva. Conseguimos entregar praticamente tudo o que planejamos, com exceção do deploy na AWS, que ficou pendente. Essa tarefa era quase exclusivamente dos AGES III, que estavam sobrecarregados com muitas tarefas e não conseguiram finalizar tudo. A reunião em si foi tranquila, e o stakeholder ficou bastante contente com os resultados que obtivemos durante a sprint.

Durante a retrospectiva da sprint, identificamos um problema recorrente em relação aos MRs, especialmente no backend, algo que pessoalmente experimentei. Notei que meus MRs ficaram pendentes por um longo período sem revisão por parte de outros colegas, o que contribuiu para minha frustração com os problemas que enfrentei. Fiquei satisfeito ao ver que nossa equipe reconheceu e abordou esse problema durante a retrospectiva. Aprendemos com essa experiência e decidimos implementar mudanças para garantir que os MRs sejam revisados de maneira mais ágil e eficiente até a próxima sprint, que será a nossa última.

Essa sprint, em geral, foi muito positiva para mim, pois aprendi várias técnicas novas ao longo dela. Desde aprofundar meu conhecimento em SpringBoot até lidar com desafios como conflitos de MRs e a implementação de deleção lógica no backend. Além disso, a colaboração com meus colegas e o suporte que recebi foram essenciais para enfrentar os desafios e alcançar os objetivos para a sprint. 


